\documentclass[10pt,a4paper]{amsart}
\usepackage[margin=19mm]{geometry}
\usepackage{amsmath,amssymb,mathtools}
\usepackage[scr=rsfs,cal=euler]{mathalfa}
\usepackage{xfrac}
\usepackage[unicode,hidelinks,bookmarks=false]{hyperref}
\newcommand{\ie}{i.e.\ }
\newcommand{\cf}{cf.\ }
\newcommand{\resp}{respectively\ }
\newcommand{\Subsection}{\subsection}
\providecommand{\nobreakdash}{\nobreak\hbox{-}}
\setlength{\emergencystretch}{2em}
\multlinegap0pt
\usepackage{bm}
\usepackage{bbm}
\usepackage{dsfont}
\usepackage{xspace}
\usepackage{cancel}
\usepackage{mathtools}
\usepackage{amsthm}
\makeatletter
\@ifundefined{theorem}{\newtheorem{theorem}{Theorem}}{}
\@ifundefined{corollary}{\newtheorem{corollary}[theorem]{Corollary}}{}
\makeatother
\newcommand{\E}{\mathcal{E}}
\newcommand{\PA}{\mathcal{P}}
\newcommand{\U}{\mathcal{U}}
\newcommand{\X}{\mathcal{X}}
\title[Lagrangian Index Policy for Restless Bandits with Average Re]{Lagrangian Index Policy for Restless Bandits with Average Reward}
\author{Konstantin Avrachenkov, Vivek S. Borkar, Pratik Shah}
\date{Source: arXiv:2412.12641v3 (CC BY 4.0)}
\begin{document}
\maketitle

\noindent\emph{Source and licence.} Lagrangian Index Policy for Restless Bandits with Average Reward. Version arXiv:2412.12641v3 (CC BY 4.0), distributed under the Creative Commons Attribution 4.0 International licence, as recorded in the arXiv licence metadata for this exact version and shown on the abstract page https://arxiv.org/abs/2412.12641v3. Full rights notice: licenses/arxiv-2412.12641-CC-BY-4.0.txt.\par\smallskip

\noindent\textit{Editorial scope.} Counted unit: the Theorem whose source label is \texttt{indep} in the article (source0.tex lines 1176--1177, source label \texttt{indep}), delivered together with the article's own complete original proof of it (source0.tex lines 1179--1210, one \texttt{proof} environment of 32 lines). No mathematics is written, repaired or re-derived: statement and proof are the article's own text line for line. Required context retained uncounted: SLLNX (lines 1085--1087). Every definition, display and label of those lines is retained unchanged. Omitted from this package and not redistributed: 1--1084, 1088--1175, 1178--1178, 1211--1535. Nothing inside a retained excerpt is omitted or altered: every retained body is delivered exactly as the article's lines are written, so no omission receipt of any kind accompanies this package. Citations: the retained excerpts carry no citation command at all, so no bibliography entry of the article is transcribed and no reference is redistributed. Reference adaptation: none was needed and none was made; every reference inside the delivered unit resolves inside it, so no unresolved reference or citation is produced. Printed slips kept literal: no printed word, symbol, hypothesis or number of the counted statement or of its proof was altered. Layout and class: the article's own class is \texttt{article} with packages including amsmath, amssymb, amsfonts, mathtools, amsthm, algorithm, algpseudocode, color, graphicx, subcaption; the wrapper is the lane's pinned \texttt{amsart} editorial wrapper at 10pt on A4, which restates the theorem-like environments the retained text needs and adds the article's own macro definitions that the retained text needs (\texttt{\textbackslash E}, \texttt{\textbackslash PA}, \texttt{\textbackslash U}, \texttt{\textbackslash X}), with extra wrapper packages (bm, bbm, dsfont, xspace, cancel, mathtools). The delivered numbering is the unit's own. Assets: no retained span contains a figure environment, a graphics-inclusion command, a PDF-page inclusion, a table or a TikZ picture; no image file is copied into this package, the package declares no asset dependency and no image file is redistributed with it. Licence: version arXiv:2412.12641v3 of this article is distributed under the Creative Commons Attribution 4.0 International licence, as recorded in the arXiv licence metadata for this exact version and shown on the abstract page https://arxiv.org/abs/2412.12641v3. A full rights notice is delivered at licenses/arxiv-2412.12641-CC-BY-4.0.txt. Characters: every retained excerpt is pure ASCII, so the delivered engine, XeLaTeX with its default Latin Modern text font, reports no missing character.
\begin{corollary}\label{SLLNX} Let $Y_n, n \geq 1,$ be integrable exchangeable random variables in $\mathbb{R}^d, d \geq 1$.  Then there exists an $\mathbb{R}^d$-valued integrable random variable $\widetilde{Y}$ measurable with respect to the associated exchangeable $\sigma$-field such that
$$\frac{1}{n}\sum_{m=1}^nY_m \to \widetilde{Y}, \ \mbox{a.s.}$$
\end{corollary}

\begin{theorem}\label{indep} The arms are asymptotically independent (as $N\uparrow\infty$) under the above Procedure 2 and the Global Attractor Hypothesis \textbf{$(\dagger)$}, i.e., for any fixed $K$ and the arms $\{X^i_n, 1 \leq i \leq K \}$, $E\left[\prod_{i=1}^Kf_i(X^i_n)\right] = \prod_{i=1}^KE\left[f_i(X^i_n)\right]$ \ $\forall \ f_i\in C_b(\X), 1 \leq i \leq K < \infty$, $\forall n$.
\end{theorem}

\begin{proof}
Recall that for each $i \geq 1$, $X^i_n, n \geq 0,$ is a controlled Markov chain controlled by the process $U^i_n, n \geq 0$.
Let $\E$ denote the $\sigma$-field of exchangeable events. Conditioned on $\E$,  $\{(X^i_n, U^i_n), n \geq 0\}, i \geq 1,$ are independent processes by de Finetti's theorem. Then almost surely,  the random probability measure $\nu_n \in \PA(\X\times\U)$ defined by
$$\nu_n(j,u) := \lim_{N\to\infty}\frac{1}{N}\sum_{i=1}^NI\{X^i_n=j, U^i_n = u\}$$
is well defined by the strong law of large numbers for exchangeable random variables, viz., Corollary \ref{SLLNX}. 
Necessarily, the random variable $\nu_n \in \PA(\X\times\U)$ is $\E$-measurable. \\

We shall write any $\kappa \in \PA(\X\times\U)$ as
$\kappa(i,u) = \bar{\kappa}(i)\varphi_\kappa(u|i) , \ i\in \X, u \in \U.$ Also define the corresponding transition probabilities as $p_\kappa(j|i) = \sum_up(j|i,u)\varphi_\kappa(u|i), \ i,j \in \X.$ In particular, disintegrating $\nu_n(x\times u)$ as $\bar{\nu}_n(x)\varphi_\nu(u|x)$, the conditional law $\varphi_\nu( \cdot | \cdot)$ is completely specified by Procedure~2 given $\bar{\nu}_n$. In particular, this allows us to write $p_{\nu_n}(j|k)$ as $p_{\bar{\nu}_n}(j|k)$ by slight abuse of notation. 
Thus, we can define
\begin{equation}
P(\bar\nu_n)= [p_{\bar\nu_n}(j|k)]_{j,k}. \label{newnotation}
\end{equation}
This notation will be important in what follows.\\

Suppose we apply Procedure 2 above. 
The procedure implies that the transition probabilities of the arms are coupled only through the dependence of the control policy on $\bar{\nu}_n$. To see this, note that a control $u \in \{0,1\}$ is assigned to $i$ by Procedure 2 depending solely on the asymptotic fraction of states (as $N\uparrow\infty$) belonging to $\gamma^{-1}(u)$ and an independent randomization if needed. Thus, the outcome of the procedure depends solely on the probability vector $\xi$ in Procedure 2, which at time $n$, corresponds to $\bar{\nu}_n$. The latter in turn is clearly $\E$-measurable. Therefore, conditioned on $\E$, for each $i$, $\{X^i_n\}$ is a Markov chain with the transition probability matrix independent of $i$ (since they are conditionally i.i.d.), but a function of $\bar{\nu}_n$ which is $\E$-measurable. \\

In addition, a reasoning analogous to the one used for defining $\nu_n$ also shows that for any $k \in \X$,
\begin{eqnarray*}
0 &=&\lim_{N\uparrow\infty}\frac{1}{N}\sum_{i=1}^N\Big(I\{X^i_{n+1} = k\} - P(X^i_{n+1} = k|X^i_n, U^i_n)\Big) \\
&=& \lim_{N\uparrow\infty}\left(\frac{1}{N}\sum_{i=1}^NI\{X^i_{n+1} = k\} - \frac{1}{N}\sum_{i=1}^N\sum_{j\in S}p_{\bar{\nu}_n}(k|j)I\{X^i_n=j\}\right) \\
&=& \bar{\nu}_{n+1}(k) - \sum_{j\in S}p_{\bar{\nu}_n}(k|j)\bar{\nu}_n(j)  \   \mbox{a.s.}, 
\end{eqnarray*}
where the first equality follows from the strong law of large numbers for martingales with bounded increments. 
Thus $\{\bar{\nu}_n\}$ satisfies the iteration
\begin{equation}
\bar{\nu}_{n+1} = \bar{\nu}_nP(\bar{\nu}_n), \ n \geq 0. \label{nuiteration}
\end{equation}

Under our \textit{Global Attractor Hypothesis} \textbf{$(\dagger)$},  $\bar{\nu}_n \to \nu^*$ a.s., where $\nu^*$ is the unique deterministic fixed point of the map $\Psi$. Hence at stationarity (i.e., as $n\to\infty$), the $\{(X^i_n, U^i_n), n \geq 0\}$ for different $i$'s are coupled only through $\nu_*$ which is deterministic. This implies their independence.
\end{proof}

\end{document}
